\begin{proof}
The normalization is finite over \(R\) and is a semilocal finite product of
Dedekind domains by \cite[Tag~0C1R]{StacksProject}; hence the displayed
branch set is finite and every \(D_{\mathfrak q}\) is a one-dimensional
normal noetherian local domain.  It is therefore a discrete valuation ring
by \cite[Tag~00PD]{StacksProject}.
Moreover, \(\mathfrak a\) is contained in no minimal prime of \(R\), while
its image on every branch lies in the maximal ideal of
\(D_{\mathfrak q}\).  Thus
\(0<d_{\mathfrak q}(\mathfrak a)<+\infty\), as asserted.

Let \(v\in\mathcal V_{\mathfrak a}(R)\), and let
\(\mathfrak p=v^{-1}(+\infty)\) be its support.  Since
\(v(\mathfrak a)=1\), the ideal \(\mathfrak a\) is not contained in
\(\mathfrak p\).  The equality
\(\sqrt{\mathfrak a}=\mathfrak m\) therefore rules out
\(\mathfrak p=\mathfrak m\).  Every nonmaximal prime of a
one-dimensional local ring is minimal, so \(\mathfrak p\) is a minimal
prime of \(R\).

The semivaluation factors through \(A=R/\mathfrak p\) and extends uniquely
to a real valuation of \(K=\operatorname{Frac}(A)\).  Let \(V\subset K\)
be its valuation ring.  A valuation ring is normal
\cite[Tag~00IC]{StacksProject}, so \(V\) contains the integral closure
\(\overline A\) of \(A\) in \(K\).  The centre
\[
                 \overline{\mathfrak q}
                  :=\mathfrak m_V\cap\overline A
\]
lies over \(\mathfrak m/\mathfrak p\).  It determines a unique
\(\mathfrak q\in\operatorname{Br}(R,\mathfrak m)\), and
\(D_{\mathfrak q}=\overline A_{\overline{\mathfrak q}}\subset V\), with
\(V\) dominating \(D_{\mathfrak q}\).

Choose a uniformizer \(\pi\) of \(D_{\mathfrak q}\).  Every nonzero
element of \(K\) has a unique expression \(u\pi^n\), where
\(u\in D_{\mathfrak q}^{\times}\) and \(n\in\mathbf Z\).  Domination gives
\(v(u)=0\) and \(c:=v(\pi)>0\), whence
\[
                              v=cw_{\mathfrak q}.
\]
Thus every real valuation occurring here has a discrete cyclic value group.
The normalization \(v(\mathfrak a)=1\) forces
\(c=d_{\mathfrak q}(\mathfrak a)^{-1}\).  Conversely, every normalized
branch order in \eqref{eq:p121-curve-branch-link} is centred at
\(\mathfrak m\) and has \(\mathfrak a\)-value one.  This proves
surjectivity and the asserted formula.

Branches with different supports are separated by their infinite loci.  If
two branches have the same support and their normalized orders agree on
\(R\), then their extensions to \(K\) are positive multiples of each other
and have the same valuation ring; their centres on \(\overline A\) are
therefore equal.  Hence the map is injective.  Pointwise convergence embeds
the finite set into the Hausdorff product
\([0,+\infty]^R\), so the resulting topology is discrete.

If \(\mathfrak b\) is \(\mathfrak m\)-primary, then every branch value
\(v(\mathfrak b)\) is finite and strictly positive.  On the branch
\(\mathfrak q\), formula \eqref{eq:p121-change-normalizing-ideal} reads
\[
 \frac{w_{\mathfrak q}}{d_{\mathfrak q}(\mathfrak a)}
 \longmapsto
 \frac{w_{\mathfrak q}}{d_{\mathfrak q}(\mathfrak b)}.
\]
The inverse and cocycle identities follow immediately.  Finally, every
excellent ring is Nagata by \cite[Tag~07QV]{StacksProject}, proving the last
assertion.
\end{proof}
